\section{Related Work}
\label{sec:related_work}

\hb{mention viescore from wenhu somewhere}

\paragraph{Video Generation Models.}
Recent advancements in video generation models have emerged from two primary architectures: diffusion-based models \cite{gen2, svd, liu2024sora, blattmann2023align, wang2023lavie, wang2023modelscope, chen2024videocrafter2, khachatryan2023text2video} and autoregressive modeling-based approaches \cite{magvit2, kondratyuk2023videopoet, hong2022cogvideo, villegas2022phenaki}. Among these, diffusion models have garnered significant attention. The model known as SVD \cite{svd}, built on a latent diffusion model (LDM) \cite{ldm}, proposes a three-stage training process for video LDMs. Given the rapid development of video generation technology, our work focuses on evaluating video generation models for their physical commonsense.

\paragraph{Evaluating Video Generation Models.} Traditional evaluation methods for video generation primarily employ metrics such as FVD \cite{unterthiner2018towards} and IS \cite{salimans2016improved}. However, there is a growing consensus on the need for more comprehensive metrics to assess the performance of video generation models \cite{huang2023vbench, liu2023evalcrafter, subjectbench, lin2024evaluating}. V-Bench \cite{huang2023vbench} offers a detailed benchmark suite that introduces a hierarchical evaluation protocol with granular perspectives. Here, we propose \name{} dataset to measure the physical commonsense in the generated videos. Additionally, we introduce a \model{} auto-evaluator and analyze specific physical laws \cite{bisk2020piqa} that are violated in the generated videos through qualitative analysis. We present a detailed related work in Appendix \ref{sec:detailed_related_work}.